Finden Sie eine Lösung der Differentialgleichungen
\begin{teilaufgaben}
\item $y'+2xy=0$
\item $(x+1)y'=3y$
\end{teilaufgaben}
mithilfe der Potenzreihenmethode.

\begin{loesung}
In beiden Teilaufgaben kann man einen Ansatz der Form
\[
y(x)
=
\sum_{n=0}^\infty a_n x^n
\qquad\text{mit}\qquad
y'(x)
=
\sum_{n=1}^\infty na_n x^{n-1}
=
\sum_{l=0}^\infty (l+1)a_{l+1} x^l
\]
benutzen.
\begin{teilaufgaben}
\item
Einsetzen in die Differentialgleichung ergibt
\begin{align*}
0
&=
\sum_{l=0}^\infty (l+1)a_{l+1} x^l
+2x
\sum_{n=0}^\infty a_n x^n
\\
&=
\sum_{l=0}^\infty (l+1)a_{l+1} x^l
+
\sum_{n=0}^\infty 2a_n x^{n+1}
\\
&=
\sum_{l=0}^\infty (l+1)a_{l+1} x^l
+
\sum_{l=1}^\infty 2a_{l-1} x^l
\\
&=
a_1
+
\sum_{l=1}^\infty \bigl((l+1)a_{l+1} + 2a_{l-1}\bigr)x^l.
\end{align*}
Koeffizientenvergleich liefert jetzt 
die Beziehungen
\begin{align*}
a_1
&=
0
\\
(l+1)a_{l+1} + 2a_{l-1}
&=
0
&&\Rightarrow&
a_{l+1} &= -\frac{2}{l+1} a_{l-1}
&&\forall l > 0.
\end{align*}
Aus der Rekursionsformel kann man bereits ablesen, dass alle
ungeraden Koeffizienten verschwinden.
Für die geraden Koeffizienten findet man iterativ
\begin{align*}
a_2 &= -a_0 \\
a_4 &= -\frac{2}{4} a_2 = \frac12 a_0 \\
a_6 &= -\frac{2}{6} a_4 = -\frac{1}{3} a_4 = -\frac{1}{3!} \\
a_8 &= -\frac{2}{8} a_6 = -\frac{1}{4} a_6 = \frac{1}{4!} \\
a_{10} &= -\frac{2}{10} a_8 = -\frac{1}{5} a_8 = -\frac{1}{5!}
 &\phantom{a}\vdots \\
a_{2k} &= (-1)^k\frac{1}{k!}.
\end{align*}
Daraus folgt die Potenzreihe
\[
y(x)
=
\sum_{n=0}^\infty
a_n x^n
=
\sum_{k=0}^\infty
a_{2k} x^{2k}
=
\sum_{k=0}^\infty
\frac{(-1)^k}{k!} (x^2)^k
=
\sum_{k=0}^\infty
\frac{1}{k!} (-x^2)^k
=
e^{-x^2}.
\]
Tatsächlich ist $y(x)=a_0e^{-x^2}$ eine Lösung der Differentialgleichung,
denn
\[
y'(x)
+
2xy(x)
=
-
a_0e^{-x^2}2x
+
2xa_0e^{-x^2}
=
0.
\]
\item
Einsetzen in die Differentialgleichung ergibt
\begin{align*}
(x+1)
\sum_{n=1}^\infty na_n x^{n-1}
&=
3
\sum_{n=0}^\infty a_n x^n
\\
\sum_{n=1}^\infty na_n x^n
+
\sum_{n=1}^\infty na_n x^{n-1}
&=
\sum_{n=0}^\infty 3a_n x^n
\\
\sum_{n=1}^\infty na_n x^n
+
\sum_{n=0}^\infty (n+1)a_{n+1} x^n
&=
\sum_{n=0}^\infty 3a_n x^n
\\
\sum_{n=1}^\infty \bigl(na_n+(n+1)a_{n+1}-3a_n\bigr) x^n
&=3a_0 - a_1.
\end{align*}
Durch Koeffizientenvergleich kann man jetzt schliessen, dass
\begin{align*}
a_1&=3a_0 \\
(n+1)a_{n+1} &= (3-n)a_n&&\forall n>0.
\end{align*}
Durch iterative Anwendung dieser Formeln findet man
\begin{align*}
a_2 &= a_1 = 3a_0\\
a_3 &= \frac{1}{3} a_2 = a_0 \\
a_4 &= \frac{0}{4} a_3 = 0.
\end{align*}
Daher ist die Lösung
\[
y(x)
=
a_0 + 3a_0 x + 3a_0 x^2 + a_0 x^3
=
a_0 (1+x)^3.
\]
Tatsächlich ist $y' = 3a_0(1+x)^2$ und daher
\[
(1+x)y'(x)
=
(1+x) 3a_0(1+x)^2
=
3a_0(1+x)^3
=
3y(x),
\]
d.~h.~die Funktion $y(x)=a_0(1+x)^3$ ist eine Lösung
der Differentialgleichung.

Die Differentialgleichung kann auch mit der Separationsmethode
direkt gelöst werden.
Dazu formt man die Gleichung erst um zu
\[
\frac{y'}{y}
=
\frac{3}{x+1}
\]
und berechnet die Stammfunktion
\begin{align*}
\int\frac{y'}{y}\,dx
&=
\int
\frac{3}{x+1}
\,dx
\intertext{auf beiden Seiten.
Die Stammfunktionen können auf beiden Seiten sofort ermittelt
werden, sie sind}
\log y &= 3\log(x+1) + C.
\intertext{Nimmet man auf beiden Seiten die Exponentialfunktion,
erhält man die Lösung}
y&=e^C(x+1)^3 = a_0(x+1)^3
\end{align*}
mit $a_0=e^C$.
\qedhere
\end{teilaufgaben}
\end{loesung}
